\chapter[1952 --- Waksman]{1952: Selman Waksman}
\label{ch:1952_selmanwaksman}

\section*{Main Contribution}

\textit{Discovered streptomycin, the first antibiotic effective against tuberculosis---a disease that had killed millions---transforming TB from a death sentence into a curable condition.}

\section{Official Citation}

for his discovery of streptomycin, the first antibiotic effective against tuberculosis

\section{Background and Historical Context}

Selman Abraham Waksman was born on July 22, 1888, in Priluki, Ukraine (then part of the Russian Empire). His family emigrated to the United States in 1910, settling in New Jersey, where Waksman would spend most of his career. He attended Rutgers University, where he earned his bachelor's degree in agriculture in 1911, and subsequently obtained his Ph.D. in soil chemistry from the University of California, Berkeley, in 1915. After completing his doctoral studies, Waksman returned to Rutgers University, where he joined the Department of Soil Microbiology and would remain for the next 40 years.

Waksman's research career was devoted to the study of soil microorganisms---the bacteria, fungi, and other microscopic organisms that inhabit the soil and play essential roles in the decomposition of organic matter, the cycling of nutrients, and the maintenance of soil health. This work, while seemingly esoteric, had significant implications for medicine, as Waksman would demonstrate that soil microorganisms are a rich source of antibiotics---chemical compounds that can kill or inhibit the growth of bacteria.

The scientific context of the 1940s was one of intense interest in the discovery of new antibiotics. Fleming's discovery of penicillin in 1928, and its subsequent development by Chain and Florey in the early 1940s, had demonstrated the enormous therapeutic potential of antibiotics. However, penicillin was effective primarily against gram-positive bacteria, and there was an urgent need for agents that could treat infections caused by gram-negative bacteria and by organisms such as \textit{Mycobacterium tuberculosis}, the bacterium that causes tuberculosis.

It was within this context of medical need that Waksman embarked on his systematic search for new antibiotics from soil microorganisms.

\section{The Discovery/Contribution}

Selman Waksman received the Nobel Prize in Physiology or Medicine ``for his discovery of streptomycin, the first antibiotic effective against tuberculosis.'' His work, carried out in collaboration with his student Albert Schatz, demonstrated that the soil bacterium \textit{Streptomyces griseus} produces an antibiotic---streptomycin---that is effective against \textit{Mycobacterium tuberculosis} and a wide range of other pathogenic bacteria.

\subsection{The Search for Antibiotics from Soil Microorganisms}

Waksman's approach to antibiotic discovery was systematic and methodical. Beginning in the 1930s, he established a large-scale screening program in which thousands of soil samples from around the world were tested for the presence of bacteria or fungi that produced antibacterial substances. The rationale behind this approach was simple but powerful: soil microorganisms exist in a highly competitive environment where they must constantly defend themselves against other microorganisms, and one of the strategies they use is the production of chemical weapons---antibiotics---that kill or inhibit their competitors.

Waksman's screening program tested soil isolates for their ability to inhibit the growth of pathogenic bacteria on agar plates. Any isolate that produced a zone of inhibition---a clear area around the colony where pathogenic bacteria could not grow---was identified as a potential antibiotic producer and subjected to further investigation. Through this systematic approach, Waksman's laboratory identified several novel antibiotics, including actinomycin (1940), streptothricin (1942), and grisein (1944).

\subsection{The Discovery of Streptomycin}

The discovery of streptomycin was the culmination of Waksman's antibiotic screening program. In 1943, Albert Schatz, a graduate student in Waksman's laboratory, isolated a bacterium from a soil sample that produced a substance capable of inhibiting the growth of \textit{Mycobacterium tuberculosis}. The bacterium was identified as \textit{Streptomyces griseus}, and the antibacterial substance it produced was named ``streptomycin.''

Streptomycin was the first antibiotic shown to be effective against tuberculosis, a disease that had been one of the leading causes of death in the world for centuries. The initial clinical trials of streptomycin, conducted in 1944 by William Feldman and Corwin Hinshaw at the Mayo Clinic, demonstrated that the drug could produce dramatic improvement in patients with pulmonary tuberculosis, including those with drug-resistant forms of the disease. These results were among a major medical breakthroughs of the twentieth century and provided the basis for the modern treatment of tuberculosis.

\subsection{The Chemistry and Mechanism of Action of Streptomycin}

Streptomycin belongs to the aminoglycoside class of antibiotics, which are characterized by their amino sugar structure and their mechanism of action. Streptomycin works by binding to the 30S subunit of the bacterial ribosome, the molecular machine responsible for protein synthesis. By binding to the ribosome, streptomycin interferes with the accuracy of the translation process, causing the ribosome to incorporate incorrect amino acids into the growing protein chain. The production of defective proteins is lethal to the bacterium, leading to cell death.

The mechanism of action of streptomycin---the disruption of protein synthesis---was fundamentally different from that of penicillin, which targets the bacterial cell wall. This difference in mechanism explained why streptomycin was effective against organisms that were resistant to penicillin, including \textit{Mycobacterium tuberculosis}. The cell wall of mycobacteria is uniquely resistant to many antibiotics, including penicillin, but streptomycin's target---the ribosome---is present in all bacteria, including mycobacteria.

\subsection{The Controversy Over Credit}

The discovery of streptomycin was the subject of a bitter dispute over credit. Albert Schatz, who had performed the initial isolation of \textit{Streptomyces griseus} and streptomycin, argued that he deserved a share of the Nobel Prize. Waksman, however, maintained that the discovery was the result of his systematic screening program and that Schatz's contribution, while important, was that of a technician carrying out a routine procedure. The dispute was resolved by a legal agreement in which Schatz received recognition for his role, but the Nobel Prize was awarded solely to Waksman.

The controversy surrounding the streptomycin discovery has been the subject of much historical analysis and debate. It highlights the complexities of assigning credit in collaborative scientific research and raises important questions about the relationship between mentor and student in the laboratory setting.

\section{Impact on Modern Medicine}

Waksman's discovery of streptomycin has had a significant impact on modern medicine, particularly in the treatment of tuberculosis and other bacterial infections.

\subsection{Treatment of Tuberculosis}

Tuberculosis (TB) was one of the leading causes of death in the world before the introduction of streptomycin. It is estimated that TB killed approximately one billion people in the twentieth century, and in many countries, TB sanatoriums---institutions where patients were isolated and given supportive care---were a common feature of the medical landscape. The introduction of streptomycin in 1944, and the subsequent development of combination chemotherapy regimens that included isoniazid (1952), pyrazinamide, and ethambutol, transformed TB from a death sentence into a curable disease.

Today, the WHO estimates that approximately 10 million people develop TB each year, and approximately 1.5 million die from the disease. While the development of drug-resistant TB strains has complicated treatment, the fundamental approach to TB chemotherapy---the use of combination drug regimens that include aminoglycoside antibiotics---remains based on the principles established by Waksman's discovery of streptomycin.

\subsection{The Antibiotic Pipeline}

Waksman's systematic approach to antibiotic discovery from soil microorganisms established a paradigm that has guided the search for new antibiotics for decades. His demonstration that soil bacteria of the genus \textit{Streptomyces} are a rich source of bioactive compounds has been confirmed many times over: approximately two-thirds of all known natural product antibiotics are produced by \textit{Streptomyces} species. The screening of soil microorganisms continues to be an important strategy for antibiotic discovery, and many clinically important antibiotics---including tetracycline, erythromycin, vancomycin, and daptomycin---were originally isolated from soil bacteria.

\subsection{Combination Chemotherapy}

Waksman's observation that streptomycin was effective against tuberculosis but that resistance could develop rapidly during monotherapy led to the development of combination chemotherapy---the use of multiple drugs simultaneously to treat infections. This approach, which is now standard of care for TB and many other bacterial infections, reduces the risk of resistance by requiring the pathogen to develop resistance to multiple drugs simultaneously. The principle of combination chemotherapy, which Waksman helped to establish, is now applied not only to infectious diseases but also to cancer treatment and other areas of medicine.

\subsection{Adverse Effects and Limitations}

Streptomycin, like all aminoglycoside antibiotics, has significant toxicities that limit its clinical use. a major of these is ototoxicity---damage to the inner ear that can cause hearing loss and balance disturbances. Streptomycin can also cause nephrotoxicity (kidney damage) and neurotoxicity. These adverse effects are particularly concerning in patients who require prolonged treatment, such as those with TB, and they have led to the development of safer alternative antibiotics for many clinical indications.

The recognition of streptomycin's toxicities has also stimulated research into the mechanisms of aminoglycoside toxicity and the development of strategies to minimize these effects. This research has been an important area of pharmacology and toxicology and has contributed to our broader understanding of drug safety.

\section{Legacy}

Selman Waksman's legacy is that of a scientist who devoted his career to the study of soil microorganisms and, in doing so, made one of the major discoveries in the history of medicine. The discovery of streptomycin saved millions of lives and opened up entirely new possibilities for the treatment of bacterial infections.

Waksman continued his research at Rutgers University after receiving the Nobel Prize, focusing on the microbiology of soil and the discovery of additional antibiotics. He founded the Institute of Microbiology at Rutgers, which became a major center for research in microbiology and antibiotic discovery. He received numerous honors during his career, including the Presidential Medal of Freedom, and died on August 16, 1973, at the age of 85.

The story of Waksman and streptomycin illustrates the power of systematic, curiosity-driven research to produce transformative medical advances. Waksman's insight that soil microorganisms might be a source of antibiotics, combined with his rigorous screening methodology, led to the discovery of a drug that influenced the trajectory of medical history. As we face the growing challenge of antimicrobial resistance, the principles and approaches that Waksman pioneered continue to be relevant and inspiring.


\section{Modern Developments}

Waksman's streptomycin discovery launched the antibiotic era for tuberculosis. Modern TB treatment uses four-drug regimens (isoniazid, rifampicin, pyrazinamide, ethambutol) for 6 months. Multidrug-resistant TB (MDR-TB) and extensively drug-resistant TB (XDR-TB) are major global threats. Bedaquiline and pretomanid, approved in 2019-2020, represent the first new TB drugs in decades. The BPaL regimen (bedaquiline, pretomanid, linezolid) cures MDR-TB in 6 months.


\section{Bibliography}

\begin{thebibliography}{9}
\bibitem{nobel-1952}
Nobel Prize Outreach, ``The Nobel Prize in Physiology or Medicine 1952,''\emph{NobelPrize.org}.\\
\url{https://www.nobelprize.org/prizes/medicine/1952/summary/}

\bibitem{lecture-1952}
Selman Waksman, ``Nobel Lecture,''\emph{NobelPrize.org}. Winner-authored primary source; downloadable from the official Nobel Prize archive.\\
\url{https://www.nobelprize.org/prizes/medicine/1952/waksman/lecture/}
\end{thebibliography}
